% Options for packages loaded elsewhere
\PassOptionsToPackage{unicode}{hyperref}
\PassOptionsToPackage{hyphens}{url}
\PassOptionsToPackage{dvipsnames,svgnames,x11names}{xcolor}
%
\documentclass[
  letterpaper,
  DIV=11,
  numbers=noendperiod]{scrreprt}

\usepackage{amsmath,amssymb}
\usepackage{iftex}
\ifPDFTeX
  \usepackage[T1]{fontenc}
  \usepackage[utf8]{inputenc}
  \usepackage{textcomp} % provide euro and other symbols
\else % if luatex or xetex
  \usepackage{unicode-math}
  \defaultfontfeatures{Scale=MatchLowercase}
  \defaultfontfeatures[\rmfamily]{Ligatures=TeX,Scale=1}
\fi
\usepackage{lmodern}
\ifPDFTeX\else  
    % xetex/luatex font selection
\fi
% Use upquote if available, for straight quotes in verbatim environments
\IfFileExists{upquote.sty}{\usepackage{upquote}}{}
\IfFileExists{microtype.sty}{% use microtype if available
  \usepackage[]{microtype}
  \UseMicrotypeSet[protrusion]{basicmath} % disable protrusion for tt fonts
}{}
\makeatletter
\@ifundefined{KOMAClassName}{% if non-KOMA class
  \IfFileExists{parskip.sty}{%
    \usepackage{parskip}
  }{% else
    \setlength{\parindent}{0pt}
    \setlength{\parskip}{6pt plus 2pt minus 1pt}}
}{% if KOMA class
  \KOMAoptions{parskip=half}}
\makeatother
\usepackage{xcolor}
\setlength{\emergencystretch}{3em} % prevent overfull lines
\setcounter{secnumdepth}{5}
% Make \paragraph and \subparagraph free-standing
\makeatletter
\ifx\paragraph\undefined\else
  \let\oldparagraph\paragraph
  \renewcommand{\paragraph}{
    \@ifstar
      \xxxParagraphStar
      \xxxParagraphNoStar
  }
  \newcommand{\xxxParagraphStar}[1]{\oldparagraph*{#1}\mbox{}}
  \newcommand{\xxxParagraphNoStar}[1]{\oldparagraph{#1}\mbox{}}
\fi
\ifx\subparagraph\undefined\else
  \let\oldsubparagraph\subparagraph
  \renewcommand{\subparagraph}{
    \@ifstar
      \xxxSubParagraphStar
      \xxxSubParagraphNoStar
  }
  \newcommand{\xxxSubParagraphStar}[1]{\oldsubparagraph*{#1}\mbox{}}
  \newcommand{\xxxSubParagraphNoStar}[1]{\oldsubparagraph{#1}\mbox{}}
\fi
\makeatother


\providecommand{\tightlist}{%
  \setlength{\itemsep}{0pt}\setlength{\parskip}{0pt}}\usepackage{longtable,booktabs,array}
\usepackage{calc} % for calculating minipage widths
% Correct order of tables after \paragraph or \subparagraph
\usepackage{etoolbox}
\makeatletter
\patchcmd\longtable{\par}{\if@noskipsec\mbox{}\fi\par}{}{}
\makeatother
% Allow footnotes in longtable head/foot
\IfFileExists{footnotehyper.sty}{\usepackage{footnotehyper}}{\usepackage{footnote}}
\makesavenoteenv{longtable}
\usepackage{graphicx}
\makeatletter
\newsavebox\pandoc@box
\newcommand*\pandocbounded[1]{% scales image to fit in text height/width
  \sbox\pandoc@box{#1}%
  \Gscale@div\@tempa{\textheight}{\dimexpr\ht\pandoc@box+\dp\pandoc@box\relax}%
  \Gscale@div\@tempb{\linewidth}{\wd\pandoc@box}%
  \ifdim\@tempb\p@<\@tempa\p@\let\@tempa\@tempb\fi% select the smaller of both
  \ifdim\@tempa\p@<\p@\scalebox{\@tempa}{\usebox\pandoc@box}%
  \else\usebox{\pandoc@box}%
  \fi%
}
% Set default figure placement to htbp
\def\fps@figure{htbp}
\makeatother
% definitions for citeproc citations
\NewDocumentCommand\citeproctext{}{}
\NewDocumentCommand\citeproc{mm}{%
  \begingroup\def\citeproctext{#2}\cite{#1}\endgroup}
\makeatletter
 % allow citations to break across lines
 \let\@cite@ofmt\@firstofone
 % avoid brackets around text for \cite:
 \def\@biblabel#1{}
 \def\@cite#1#2{{#1\if@tempswa , #2\fi}}
\makeatother
\newlength{\cslhangindent}
\setlength{\cslhangindent}{1.5em}
\newlength{\csllabelwidth}
\setlength{\csllabelwidth}{3em}
\newenvironment{CSLReferences}[2] % #1 hanging-indent, #2 entry-spacing
 {\begin{list}{}{%
  \setlength{\itemindent}{0pt}
  \setlength{\leftmargin}{0pt}
  \setlength{\parsep}{0pt}
  % turn on hanging indent if param 1 is 1
  \ifodd #1
   \setlength{\leftmargin}{\cslhangindent}
   \setlength{\itemindent}{-1\cslhangindent}
  \fi
  % set entry spacing
  \setlength{\itemsep}{#2\baselineskip}}}
 {\end{list}}
\usepackage{calc}
\newcommand{\CSLBlock}[1]{\hfill\break\parbox[t]{\linewidth}{\strut\ignorespaces#1\strut}}
\newcommand{\CSLLeftMargin}[1]{\parbox[t]{\csllabelwidth}{\strut#1\strut}}
\newcommand{\CSLRightInline}[1]{\parbox[t]{\linewidth - \csllabelwidth}{\strut#1\strut}}
\newcommand{\CSLIndent}[1]{\hspace{\cslhangindent}#1}

\KOMAoption{captions}{tableheading}
\makeatletter
\@ifpackageloaded{bookmark}{}{\usepackage{bookmark}}
\makeatother
\makeatletter
\@ifpackageloaded{caption}{}{\usepackage{caption}}
\AtBeginDocument{%
\ifdefined\contentsname
  \renewcommand*\contentsname{Table of contents}
\else
  \newcommand\contentsname{Table of contents}
\fi
\ifdefined\listfigurename
  \renewcommand*\listfigurename{List of Figures}
\else
  \newcommand\listfigurename{List of Figures}
\fi
\ifdefined\listtablename
  \renewcommand*\listtablename{List of Tables}
\else
  \newcommand\listtablename{List of Tables}
\fi
\ifdefined\figurename
  \renewcommand*\figurename{Figure}
\else
  \newcommand\figurename{Figure}
\fi
\ifdefined\tablename
  \renewcommand*\tablename{Table}
\else
  \newcommand\tablename{Table}
\fi
}
\@ifpackageloaded{float}{}{\usepackage{float}}
\floatstyle{ruled}
\@ifundefined{c@chapter}{\newfloat{codelisting}{h}{lop}}{\newfloat{codelisting}{h}{lop}[chapter]}
\floatname{codelisting}{Listing}
\newcommand*\listoflistings{\listof{codelisting}{List of Listings}}
\makeatother
\makeatletter
\makeatother
\makeatletter
\@ifpackageloaded{caption}{}{\usepackage{caption}}
\@ifpackageloaded{subcaption}{}{\usepackage{subcaption}}
\makeatother

\usepackage{bookmark}

\IfFileExists{xurl.sty}{\usepackage{xurl}}{} % add URL line breaks if available
\urlstyle{same} % disable monospaced font for URLs
\hypersetup{
  pdftitle={Workplace Injury Risk: Office Workers vs.~Manufacturing Workers},
  pdfauthor={Brynn},
  colorlinks=true,
  linkcolor={blue},
  filecolor={Maroon},
  citecolor={Blue},
  urlcolor={Blue},
  pdfcreator={LaTeX via pandoc}}


\title{Workplace Injury Risk: Office Workers vs.~Manufacturing Workers}
\author{Brynn}
\date{2026-04-20}

\begin{document}
\maketitle

\renewcommand*\contentsname{Table of contents}
{
\hypersetup{linkcolor=}
\setcounter{tocdepth}{2}
\tableofcontents
}

\bookmarksetup{startatroot}

\chapter{Executive Summary}\label{executive-summary}

\section{Executive Summary}\label{executive-summary-1}

This report compares workplace injury risk between manufacturing workers
and office workers (finance, insurance, and professional services) using
2023--2024 data from the U.S. Bureau of Labor Statistics. The analysis
covers both nonfatal injuries and workplace deaths.

Key findings:

\begin{itemize}
\tightlist
\item
  \textbf{Manufacturing workers are injured at a much higher rate
  overall} --- 86.4 injuries per 10,000 workers, compared to 11.0 in
  finance and insurance and 18.6 in professional services.
\item
  \textbf{Falls make up a disproportionately large share of injuries in
  office sectors.} Falls account for 33.6\% of all injuries in finance
  and insurance, compared to 21.6\% in manufacturing.
\item
  \textbf{The same pattern holds for workplace deaths.} Falls caused
  33.3\% of deaths in finance and insurance, versus 14.7\% in
  manufacturing.
\item
  \textbf{Recommendation:} Office environments should prioritize fall
  prevention programs even though overall injury rates are low. Falls
  are the dominant injury risk in these settings.
\end{itemize}

\begin{center}\rule{0.5\linewidth}{0.5pt}\end{center}

\bookmarksetup{startatroot}

\chapter{Introduction}\label{introduction}

\section{Introduction}\label{introduction-1}

Most people assume that factories are far more dangerous than offices,
and when it comes to overall injury rates, that assumption is correct.
Manufacturing involves heavy machinery, physical labor, and industrial
hazards that drive injury rates well above most other sectors.

But one specific hazard tells a more complicated story: falls. Slippery
floors, cluttered walkways, and everyday movement in office settings can
lead to falls that cause serious injuries. Understanding how fall risk
compares across sectors matters for safety officers who need to allocate
resources effectively.

This report addresses the question: \emph{Are office workers at greater
risk of falls than manufacturing workers, and how do injuries compare in
severity?} To answer this, we use two federal datasets from the Bureau
of Labor Statistics (BLS) that together cover the full spectrum of
workplace injury, from minor incidents requiring days off work to fatal
injuries.

\begin{center}\rule{0.5\linewidth}{0.5pt}\end{center}

\bookmarksetup{startatroot}

\chapter{Data}\label{data}

\section{Data}\label{data-1}

\subsection{Nonfatal Injuries: SOII}\label{nonfatal-injuries-soii}

The Survey of Occupational Injuries and Illnesses (SOII) is published
annually by the BLS, as required by the Occupational Safety and Health
Act of 1970 (Bureau of Labor Statistics n.d.b). It surveys approximately
200,000 employers and reports injury rates per 10,000 full-time
equivalent (FTE) workers. The most recent release covers 2023--2024
(Bureau of Labor Statistics 2026b).

For this report, we used Table R8, which breaks down nonfatal injury
rates by industry and event type (including falls, slips, and trips). We
focused on Days Away From Work (DAFW) rates, which capture injuries
serious enough to keep workers out of work.

\subsection{Fatal Injuries: CFOI}\label{fatal-injuries-cfoi}

The Census of Fatal Occupational Injuries (CFOI) is an annual count of
every workplace death in the United States (Bureau of Labor Statistics
n.d.a). Unlike the SOII, which is a sample survey, the CFOI verifies
each case using multiple independent sources such as death certificates,
OSHA reports, and workers' compensation records. The most recent release
covers 2024 (Bureau of Labor Statistics 2026a).

We used Table A-1, which reports fatal injury counts by industry and
event type, including falls.

\bookmarksetup{startatroot}

\chapter{Methods}\label{methods}

Injury rates for three sectors were extracted from the SOII Table R8
using NAICS codes: manufacturing (31--33), finance and insurance (52),
and professional, scientific, and technical services (54). Finance and
insurance and professional services were selected as the most
representative office-type sectors.

For each sector, the following are recorded:

\begin{itemize}
\tightlist
\item
  \textbf{Total DAFW rate} --- overall nonfatal injury rate per 10,000
  FTE workers
\item
  \textbf{Fall-specific DAFW rate} --- nonfatal fall injuries per 10,000
  FTE workers
\item
  \textbf{Fall share of total injuries} --- falls as a percentage of all
  nonfatal injuries
\end{itemize}

Fatal injury counts were taken directly from CFOI Table A-1. Fall share
of deaths are calculated by dividing fatal falls by total sector deaths.

No statistical tests were conducted. Both the SOII and CFOI represent
population-level estimates, so the observed differences between sectors
reflect real differences rather than sampling error.

\bookmarksetup{startatroot}

\chapter{Results}\label{results}

\section{Results}\label{results-1}

\subsection{Nonfatal Injury Rates}\label{nonfatal-injury-rates}

As shown in Table~\ref{tbl-nonfatal}, manufacturing workers experienced
a total nonfatal injury rate of 86.4 per 10,000 FTE workers --- roughly
8 times higher than the rate in finance and insurance (11) and 4.6 times
higher than professional services (18.6).

\begin{longtable}[]{@{}lrrl@{}}

\caption{\label{tbl-nonfatal}Nonfatal injury rates per 10,000 FTE
workers by sector (SOII, 2023--2024)}

\tabularnewline

\toprule\noalign{}
Sector & Total Rate & Fall Rate & Falls as \% of Total \\
\midrule\noalign{}
\endhead
\bottomrule\noalign{}
\endlastfoot
Manufacturing & 86.4 & 18.7 & 21.6\% \\
Finance \& Insurance & 11.0 & 3.7 & 33.6\% \\
Professional Services & 18.6 & 4.3 & 23.1\% \\

\end{longtable}

\subsection{Fall-Specific Risk}\label{fall-specific-risk}

While manufacturing workers have more falls in absolute terms, falls
represent a much larger \emph{proportion} of injuries in office sectors.
As seen in Table~\ref{tbl-nonfatal}, falls account for 33.6\% of all
nonfatal injuries in finance and insurance, and 23.1\% in professional
services --- compared to just 21.6\% in manufacturing.

This means that in office environments, falls are the dominant injury
type. In manufacturing, falls are just one of many hazards competing for
attention.

\subsection{Fatal Injuries}\label{fatal-injuries}

Table~\ref{tbl-fatal} shows fatal injury counts and fall-related deaths
by sector in 2024. Manufacturing had 353 workplace deaths, with 52
(14.7\%) caused by falls. Finance and insurance had 18 total deaths,
with 6 (33.3\%) from falls. Professional services had 65 deaths, with 10
(15.4\%) from falls.

\begin{longtable}[]{@{}
  >{\raggedright\arraybackslash}p{(\linewidth - 6\tabcolsep) * \real{0.2973}}
  >{\raggedleft\arraybackslash}p{(\linewidth - 6\tabcolsep) * \real{0.1757}}
  >{\raggedleft\arraybackslash}p{(\linewidth - 6\tabcolsep) * \real{0.2432}}
  >{\raggedright\arraybackslash}p{(\linewidth - 6\tabcolsep) * \real{0.2838}}@{}}

\caption{\label{tbl-fatal}Fatal occupational injuries by sector (CFOI,
2024)}

\tabularnewline

\toprule\noalign{}
\begin{minipage}[b]{\linewidth}\raggedright
Sector
\end{minipage} & \begin{minipage}[b]{\linewidth}\raggedleft
Total Deaths
\end{minipage} & \begin{minipage}[b]{\linewidth}\raggedleft
Deaths from Falls
\end{minipage} & \begin{minipage}[b]{\linewidth}\raggedright
Falls as \% of Deaths
\end{minipage} \\
\midrule\noalign{}
\endhead
\bottomrule\noalign{}
\endlastfoot
Manufacturing & 353 & 52 & 14.7\% \\
Finance \& Insurance & 18 & 6 & 33.3\% \\
Professional Services & 65 & 10 & 15.4\% \\

\end{longtable}

The pattern is consistent across both nonfatal and fatal data: falls are
proportionally far more significant in office sectors than in
manufacturing.

\subsection{Summary of Key
Comparisons}\label{summary-of-key-comparisons}

Table~\ref{tbl-summary} provides a side-by-side comparison of the fall
share of injuries across both nonfatal and fatal measures.

\begin{longtable}[]{@{}
  >{\raggedright\arraybackslash}p{(\linewidth - 4\tabcolsep) * \real{0.3099}}
  >{\raggedright\arraybackslash}p{(\linewidth - 4\tabcolsep) * \real{0.4225}}
  >{\raggedright\arraybackslash}p{(\linewidth - 4\tabcolsep) * \real{0.2676}}@{}}

\caption{\label{tbl-summary}Falls as a share of total injuries ---
nonfatal and fatal comparison}

\tabularnewline

\toprule\noalign{}
\begin{minipage}[b]{\linewidth}\raggedright
Sector
\end{minipage} & \begin{minipage}[b]{\linewidth}\raggedright
Falls: \% of Nonfatal Injuries
\end{minipage} & \begin{minipage}[b]{\linewidth}\raggedright
Falls: \% of Deaths
\end{minipage} \\
\midrule\noalign{}
\endhead
\bottomrule\noalign{}
\endlastfoot
Manufacturing & 21.6\% & 14.7\% \\
Finance \& Insurance & 33.6\% & 33.3\% \\
Professional Services & 23.1\% & 15.4\% \\

\end{longtable}

\bookmarksetup{startatroot}

\chapter{Conclusion}\label{conclusion}

\section{Conclusion}\label{conclusion-1}

This report looked at the workplace injury risk across manufacturing and
office sectors using both BLS nonfatal and fatal injury data. The
findings confirm that manufacturing workers face much higher injury
rates overall, but that falls are disproportionately common in office
environments.

The nonfatal injury data showed that falls account for 33.6\% of all
injuries in finance and insurance, compared to just 21.6\% in
manufacturing. The fatal injury data showed the same pattern: falls
caused 33.3\% of deaths in finance and insurance versus 14.7\% in
manufacturing. T

These results suggest clear and practical next steps: safety officers in
office settings should treat fall prevention as a top priority. Even
though overall injury numbers are low, when office workers do get hurt,
it is most likely from a fall. Simple changes such as slip-resistant
flooring, clear walkways, proper lighting, and employee awareness
programs can help to reduce this risk.

There are some limitations to note. The SOII relies on employer
self-reporting, which may undercount injuries in workplaces with weaker
safety cultures. Broad NAICS categories also include a wide range of job
types, so rates may not perfectly represent desk-based workers
specifically. Median days away from work (a useful measure of injury
severity) was not available in the tables used here and would strengthen
future analyses.

Overall, the data shows a clear case for workplace safety, recognizing
fall risk as a significant concern even in low-injury office
environments.

\phantomsection\label{refs}
\begin{CSLReferences}{1}{0}
\bibitem[\citeproctext]{ref-bls_cfoi_2024}
Bureau of Labor Statistics. 2026a. {``Census of Fatal Occupational
Injuries Summary, 2024.''}
https://www.bls.gov/news.release/cfoi.nr0.htm.

\bibitem[\citeproctext]{ref-bls_soii_2024}
---------. 2026b. {``Employer-Reported Workplace Injuries and Illnesses,
2023--2024.''} https://www.bls.gov/news.release/osh.nr0.htm.

\bibitem[\citeproctext]{ref-bls_cfoi_tables}
---------. n.d.a. {``Census of Fatal Occupational Injuries ({CFOI}) ---
Current and Revised Data.''}
https://www.bls.gov/iif/fatal-injuries-tables.htm. Accessed February 21,
2026.

\bibitem[\citeproctext]{ref-bls_soii_tables}
---------. n.d.b. {``Survey of Occupational Injuries and Illnesses
Data.''}
https://www.bls.gov/iif/nonfatal-injuries-and-illnesses-tables.htm.
Accessed February 21, 2026.

\end{CSLReferences}




\end{document}
